\section{数据到文本}
\label{sec:nlp-generation-data-to-text}

\term{数据到文本生成}（data-to-text generation）以表格、记录或三元组为输入，生成描述它们的自然语言。输入已经是结构化的事实，模型不必再识别信息，难点转到表达上：在许多记录中选出该说的几条，安排叙述顺序，把数值和名称准确写出，并且不补造记录中没有的事实。这四项分别对应第~\ref{sec:nlp-generation-challenges}小节的内容选择、重排、复制和忠实性挑战。

方法谱系沿这四项展开。人工流水线把生成拆成若干阶段，每一阶段由规则完成，可控、可检查，但规则难以迁移到新领域。端到端神经生成去掉人工规则，输出流畅，却常写出记录不支持的数字和陈述，内容选择也不稳定。显式内容选择与规划把流水线的前两个阶段改成可学习的模块，在写句子之前先产生可检查的计划。结构线性化后交给预训练编码器—解码器或大语言模型，则借用预训练学到的语言能力改善流畅性和低资源表现。节末表~\ref{tab:nlp-generation-d2t-methods}比较四者。

本节使用五条规范化记录，每条写成（实体，类型，值）三元组，实体均为活动编号\code{ML-0418}：$r_1$的类型为名称，值为“机器学习读书会”；$r_2$为时间“2025年4月18日14:00”；$r_3$为地点“海棠楼201”；$r_4$为校内规则“校内师生可直接参加”；$r_5$为校外规则“校外来宾须在2025年4月16日18:00前提交报名信息”。这是从原通知抽出的当前安排视图，没有原日期字段，因而不能单凭这张表生成“从4月17日改期”的历史陈述。

\subsection{经典生成流水线}
\label{sec:nlp-generation-classic-pipeline}

在神经方法之前，实用的生成系统多采用Reiter与Dale总结的流水线\scite{nlp-generation-reiter2000building}。它把“从数据到句子”拆成三个模块：\term{文档规划}（document planning）决定说什么、按什么结构说，包括内容确定与文档结构化两项任务；\term{微规划}（microplanning）决定用什么词、怎样指代、哪些事实合并成一句；\term{表层实现}（surface realization）按语法把句子计划变成字符串，处理词序、虚词和标点。每一阶段的输出都是显式的中间表示，可以单独检查和修改，这是流水线的主要优点；代价是规则和模板需要按领域人工编写，换领域或换语言时重写成本高，表达也较单一。

\subsubsection{内容确定、文档规划与微规划}
\label{sec:nlp-generation-nlg-pipeline-stages}

\term{内容确定}（content determination）从全部记录中选出要报道的条目，通常由领域规则完成。设任务为“给校外来宾的两句通知”，规则可以写成：名称、时间、地点必选；校内规则与读者无关，不选；校外规则必选。五条记录因此选出$r_1,r_2,r_3,r_5$，得到四条待表达的消息。

文档结构化把消息组织成有序的文档计划，常用修辞关系连接。这里第一句介绍活动本身（$r_1,r_2,r_3$），第二句说明参加条件（$r_5$），两句之间是“说明”关系：后者补充前者的参加方式。文档计划可以写成树：根节点为“说明”，左子树是由$r_1,r_2,r_3$组成的事件描述，右子树是$r_5$。

微规划把文档计划变成句子计划。词汇化为每条消息选词，例如时间类型选“于……举行”而不是“定在”；指代表达生成决定第二次提到活动时用全名、简称还是省略，第二句以“校外来宾”开头，不必再提活动名；句子聚合把同一事件的名称、时间和地点合并成一个小句，而不是写成三句“活动名为……。时间为……。地点为……。”。最后，表层实现按中文语序输出
\begin{quote}
机器学习读书会于4月18日14:00在海棠楼201举行。校外来宾须在4月16日18:00前提交报名信息。
\end{quote}
日期按规则从完整格式缩写为“4月18日”，这一步也要记录在实现规则中，否则读者无法确认年份是否被有意省略。

\subsubsection{模板与槽位填充}
\label{sec:nlp-generation-template}

字段和句式稳定时，可以把微规划与表层实现合并成模板：
“\{名称\}于\{时间\}在\{地点\}举行。\{校外规则\}。”
模板不需要训练，参数来自经确认的模式、日期格式和语法规则。执行时先按活动编号连接字段，统一时区与日期，再填槽；不能把两条不同活动的同名字段按表格行号直接拼接。

填入上述记录后，第一句的三个槽位分别回指$r_1,r_2,r_3$，第二句回指$r_5$，每个陈述都能追溯到一条记录。若地点缺失，模板应选择明确的缺值分支，如“地点尚未提供”，或者停止生成并报告缺字段，不能填入默认教室。两条完全相同的地点记录只表达一次，但日期相同而编号不同的活动不能去重：去重按对象身份进行，不按文字进行。

模板的事实覆盖可控，弱点在组合。每增加一种字段组合、一种缺值情形或一种语气，模板数就要成倍增加；大量记录写成长文时，固定句式反复出现，读起来机械。把同样的记录直接交给条件生成器，表达更灵活，却可能写出“读书会周五下午举行，欢迎参加”：语句通顺，却没有表达精确时间、地点和报名规则，以五条记录为单位核对，它至多完整覆盖$r_1$。下一小节的端到端方法正是从这种取舍出发的。

\subsection{端到端神经生成与无依据内容}
\label{sec:nlp-generation-end-to-end-d2t}

端到端方法的动机是去掉人工规则：把全部记录交给编码器—解码器，由数据自动学出选择、排序和措辞。Wiseman等在体育比赛数据上系统地检验了这一路线\scite{nlp-generation-wiseman2017}。他们构建的RotoWire数据集把NBA比赛的球员与球队统计表与人写的赛后报道配对，每场比赛平均约有六百条记录，而一篇报道只提到其中一小部分，因此内容选择是主要难点。

模型的输入是记录集合$R=\{r_1,\ldots,r_J\}$，每条记录有实体$e_j$、类型$t_j$和值$m_j$（原数据还有主客队标记）。记录编码把各字段嵌入拼接后经一层前馈网络：
\[
  \bm u_j=\bigl[\bm E_e(e_j);\bm E_t(t_j);\bm E_m(m_j)\bigr],\qquad
  \bm r_j=\operatorname{ReLU}(\bm W_r\bm u_j+\bm b_r)\in\mathbb R^{d}.
\]
解码器是带注意力的LSTM，第$t$步按式~\eqref{eq:nlp-generation-bahdanau}对全部$\bm r_j$求注意力$a_{t,j}$。复制单位是记录的值：门控复制按
\[
  p_t(w)=p(z_t=1\mid\bm s_t)\sum_{j:\,m_j=w}a_{t,j}+p(z_t=0\mid\bm s_t)\,P_{\mathrm{vocab}}(w)
\]
混合，与式~\eqref{eq:nlp-generation-pointer-generator}相同，只是复制分布定义在记录值而不是源端词元上；他们也比较了CopyNet式的共享归一化复制。为了让解码器状态覆盖更多输入记录，他们还试用了只在训练时使用的重建损失：把解码器隐状态切成长度不超过$B$的连续块，要求从每一块预测出表中若干条记录的实体、类型与值，以表中最匹配记录的负对数似然计损失，作为辅助目标加到式~\eqref{eq:nlp-generation-nll}上。

只看BLEU无法判断报道说对了什么，Wiseman等因此提出\term{抽取式评价}（extractive evaluation）：先在参考报道上训练一个关系抽取器，再用它从生成文本$\hat y$中抽出（实体，值，类型）关系序列$E(\hat y)$，从参考文本中抽出$E(y^{\star})$，计算三项指标。关系生成（RG）统计$E(\hat y)$中有多少条能在记录$R$中找到，衡量忠实性；内容选择（CS）以$E(y^{\star})$为参照计算$E(\hat y)$的精确率与召回率，衡量是否说了人会说的事；内容排序（CO）用两条关系序列之间的归一化Damerau--Levenshtein距离衡量顺序是否相近。例如生成句“读书会于4月18日14:00在海棠楼301举行，主讲人为张老师”抽出三条关系，只有时间一条在记录中，RG精确率为$1/3$；地点值写错，主讲人是记录没有的补造内容（教学示例）。

评价结果揭示了端到端路线的典型问题。神经模型的输出比模板流畅得多，但RG精确率明显低于模板：模型会写出记录中不存在的数字、把一名球员的数据安到另一名球员身上，或补写记录无法支持的评论；CS也不稳定，同一模型时而遗漏关键记录、时而罗列次要统计。原因在于所有决定都隐含在逐词生成里：注意力每步只局部地决定读哪条记录，没有一个全局计划约束“这篇报道要讲哪几条、按什么顺序”。抽取式评价本身也依赖抽取器的精度，抽取错误会同时影响三项指标。下一小节把内容选择与规划重新显式化。

\subsection{内容选择与规划后生成}
\label{sec:nlp-generation-content-planning}

Puduppully等的方法把流水线中的内容确定与文档规划变成可学习的模块，并保留端到端训练\scite{nlp-generation-puduppully2019}。\term{内容选择与规划}（content selection and planning）先从记录集合$R$产生记录编号序列$z=(z_1,\ldots,z_K)$作为内容计划，再以计划为条件写文本。图~\ref{fig:nlp-generation-content-planning}给出五个计算阶段。

\begin{figure*}[tbp]
  \centering
  % 内容选择与规划后生成（Puduppully等，2019）的五个计算阶段及两项训练损失。
% 下方示例行使用正文的五条活动记录，为教学示例。
\begin{tikzpicture}[
  x=1mm,y=1mm,figure picture,
  stage/.style={figure node,minimum width=24mm,text width=22mm,minimum height=17mm,
    inner sep=2pt,font=\figurefont\scriptsize,align=flush center},
  ex/.style={font=\figurefont\scriptsize,text=FigureMuted,text width=26mm,align=flush center},
  loss/.style={figure node,draw=FigureLoss,fill=FigureLossFill,minimum height=6mm,
    inner sep=1.5pt,font=\figurefont\scriptsize},
  flow/.style={figure flow}
]
  \node[stage,draw=FigureInput,fill=FigureInputFill] (enc) at (0,0)
    {记录编码\\字段嵌入拼接后\\经前馈层得$\bm r_j$};
  \node[stage,draw=FigureModel] (gate) at (30,0)
    {内容选择门\\$\bm r^{\mathrm{cs}}_j=\bm g_j\odot\bm r_j$};
  \node[stage,draw=FigureModel] (plan) at (60,0)
    {指针规划\\$p(z_k=j\mid z_{<k},R)$};
  \node[stage,draw=FigureModel] (penc) at (90,0)
    {计划编码\\双向LSTM得$\bm e_k$};
  \node[stage,draw=FigureOutput,fill=FigureOutputFill] (dec) at (120,0)
    {文本解码\\注意力$\bm e_k$与复制};
  \foreach \a/\b in {enc/gate,gate/plan,plan/penc,penc/dec}{
    \draw[flow] (\a.east) -- (\b.west);
  }
  \node[ex,anchor=north] at (enc.south) {$r_1,\ldots,r_5$各为\\（实体，类型，值）};
  \node[ex,anchor=north] at (gate.south) {每条记录读取\\其他记录后得门值};
  \node[ex,anchor=north] at (plan.south) {$z=(r_1,r_2,r_3,r_5,$\\$\mathrm{EOS})$};
  \node[ex,anchor=north] at (penc.south) {只编码计划中的\\四条记录};
  \node[ex,anchor=north] at (dec.south) {“读书会于4月18日\\14:00在海棠楼201……”};
  \node[loss] (lp) at (60,14) {$-\log p_{\bm\phi}(z\mid R)$};
  \node[loss] (ly) at (120,14) {$-\log p_{\bm\theta}(y\mid R,z)$};
  \draw[figure link,densely dashed] (lp) -- (plan.north);
  \draw[figure link,densely dashed] (ly) -- (dec.north);
\end{tikzpicture}

  \caption{内容选择与规划后生成的计算阶段。记录编码与端到端方法相同；内容选择门让每条记录参照其他记录调节自身表示；指针网络在记录集合上逐步选出内容计划，并以结束符终止；计划经双向LSTM编码后，由带注意力与复制的解码器写出文本。虚线框为两项训练损失，分别监督计划与文本。下方示例行使用本节的五条活动记录，为教学示例。}
  \label{fig:nlp-generation-content-planning}
\end{figure*}

第一步是内容选择门。一条记录是否值得报道，取决于其他记录：一名球员得分很高才值得提，平常的数值则不必。门控让每条记录先注意其他记录，再据此缩放自身表示：
\begin{equation}
  \begin{gathered}
  \alpha_{j,k}=\frac{\exp(\bm r_j^{\top}\bm W_a\bm r_k)}{\sum_{k'\neq j}\exp(\bm r_j^{\top}\bm W_a\bm r_{k'})},\qquad
  \bm c_j=\sum_{k\neq j}\alpha_{j,k}\bm r_k,\\
  \bm g_j=\sigma\bigl(\bm W_g[\bm r_j;\bm c_j]\bigr),\qquad
  \bm r^{\mathrm{cs}}_j=\bm g_j\odot\bm r_j .
  \end{gathered}
  \label{eq:nlp-generation-content-gate}
\end{equation}
$\bm c_j$汇总与第$j$条记录相关的其他记录，门向量$\bm g_j\in(0,1)^d$逐维决定保留多少，$\odot$是逐元素乘法。门没有单独的标签，随后续两项损失一起训练。

第二步是指针规划。一个LSTM解码器在记录集合上运行指针网络：初始状态取全部$\bm r^{\mathrm{cs}}_j$的平均，第$k$步以上一选中记录的$\bm r^{\mathrm{cs}}$为输入得到状态$\bm h_k$，按$p(z_k=j\mid z_{<k},R)\propto\exp(\bm h_k^{\top}\bm W_c\bm r^{\mathrm{cs}}_j)$在记录上归一化。它与式~\eqref{eq:nlp-generation-pointer-network}的区别只在打分用双线性形式而非加性形式，输出同样是输入位置。计划长度不固定，实现时在记录集合中加一个特殊的结束记录，选中它即停止。第三步，把选中的记录$\bm r^{\mathrm{cs}}_{z_1},\ldots,\bm r^{\mathrm{cs}}_{z_K}$依次送入双向LSTM，得到计划编码$\bm e_1,\ldots,\bm e_K$。第四步，文本解码器只对$\bm e_k$做注意力，并从计划中的记录复制值。解码器不再面对数百条记录，只需把一条短计划依次写出来。

训练需要记录、计划和文本三种监督，目标是两项负对数似然之和：
\begin{equation}
  L_{\mathrm{plan}}
  =-\log p_{\bm{\phi}}(z\mid R)
   -\log p_{\bm{\theta}}(y\mid R,z).
  \label{eq:nlp-generation-plan-loss}
\end{equation}
第一项监督选择与次序，第二项监督如何把选定记录说出来。数据集只有记录与文本，计划标签要另行构造：用上一小节的关系抽取器从参考文本中抽出关系，按出现顺序匹配到记录，得到弱监督的计划。推断时两个阶段各用一次束搜索：先求近似最优的计划$\hat z$，再以$\hat z$为条件生成文本。

以给校外来宾的两句通知为例，计划为$z=(r_1,r_2,r_3,r_5,\mathrm{EOS})$。若规划器给这五步的概率依次为$0.7,0.6,0.8,0.9,0.95$，第一项为$-\log(0.7\times0.6\times0.8\times0.9\times0.95)=-\log0.28728\approx1.2473$；文本部分沿用第~\ref{sec:nlp-generation-training}小节的三词元算例，第二项为$1.4271$，总损失约$2.6744$（教学数值）。两项可以分别监控：计划损失下降而文本损失不降，说明问题在表达；反之说明选错了内容。

显式计划带来两项好处。其一，内容选择可以在写句子之前单独检查：计划中有$r_5$而无$r_4$，正符合来宾通知的用途；若任务改成完整通知，同样的计划就不合格。其二，解码器只读计划中的少数记录，补造其他记录内容的机会减少。在RotoWire测试集上，Puduppully等报告的RG精确率从Wiseman等端到端基线的约75\%提高到约87\%。代价同样来自计划：弱对齐的计划标签含有抽取错误；推断要多运行一次规划解码；计划一旦漏选某条记录，文本解码器只读计划，无法再补回。复制能写出“海棠楼201”，但不阻止模型把它绑定到另一场活动，生成后仍要把每个陈述对齐回$r_j$核验。

\subsection{结构线性化与预训练模型}
\label{sec:nlp-generation-linearization}

前两小节的模型从零训练，数据到文本的平行语料又通常只有数千到数万对，模型要同时学会语言和数据，表达能力受限。预训练编码器—解码器已经掌握了语言表达，剩下的问题是把结构送进一个只接受词元序列的模型。最常用的做法是\term{线性化}（linearization）：按固定顺序把结构展开成带分隔标记的文本。三元组可以写成
\begin{center}
  \code{<H> ML-0418 <R> 时间 <T> 2025-04-18 14:00 <H> ML-0418 <R> 地点 <T> 海棠楼201}
\end{center}
其中\code{<H>}、\code{<R>}、\code{<T>}分别标出主语、关系和宾语；表格可以按行展开，每个单元格前加列名与行名。线性化后的序列作为编码器输入，参考文本作为目标，按式~\eqref{eq:nlp-generation-nll}微调T5或BART；两者的去噪预训练目标见第~\ref{sec:nlp-llm-objective-t5-bart}小节。Kale与Rastogi用这一做法微调T5，在表格到文本、三元组到文本和对话行为到文本等任务上都超过了专门设计的结构编码器，在训练中未见过的类别上提高尤其明显\scite{nlp-generation-kale2020t5d2t}。

两个数据集体现了这一阶段的任务设置。WebNLG以DBpedia中的一到七条RDF三元组为输入，测试集包含训练时未出现的类别，用来检验对新领域关系的泛化\scite{nlp-intro-gardent2017}。ToTTo以维基百科表格加一组高亮单元格为输入，目标是一句只描述高亮内容的句子\scite{nlp-generation-parikh2020totto}。高亮单元格把内容选择从模型手中拿走，任务集中在忠实表达上；参考句由标注者从原维基句子逐步修改而来，删去表格不支持的内容，因此适合衡量忠实性。

线性化把结构信息变成了顺序信息，带来三个实现问题。第一，顺序本身没有含义，但模型会对顺序敏感：同一组三元组换一种排列，输出可能不同。训练时随机打乱三元组顺序可以减少这种依赖。第二，长表格超出编码器长度时只能截断或先选子表，ToTTo的高亮单元格正是一种预先选择。第三，数值和日期按子词切分后，“2025-04-18”可能被拆成若干片段，格式不统一时模型容易抄错；应统一数值格式，必要时让模型输出记录编号或占位符，再由程序填值。

大语言模型可以直接据提示完成数据到文本：把记录或表格放进提示，说明读者与篇幅，必要时给几个示例。这省去了微调数据，却把内容选择重新变成隐式的：模型自己决定说哪些记录，也可能用常识补写记录没有的内容。对大模型输出，应把每个数值和实体回查到记录字段，按第~\ref{sec:nlp-generation-recovery}小节的流程核验；需要可审计的内容选择时，可以先让模型输出记录编号形式的计划，经检查后再写文本，这正是上一小节显式规划在提示层面的对应做法。

表~\ref{tab:nlp-generation-d2t-methods}比较本节四类方法。预训练路线显著改善了流畅性和低资源表现，显式规划仍在可审计性上有优势，两者可以结合：规划部分输出记录编号，表达部分由预训练模型完成。

\begin{table*}[htbp]
  \centering
  \small
  \begin{tabularx}{\linewidth}{@{}>{\raggedright\arraybackslash}p{24mm}
    >{\raggedright\arraybackslash}X>{\raggedright\arraybackslash}X
    >{\raggedright\arraybackslash}X@{}}
    \toprule
    方法 & 解决的问题 & 核心机制 & 新代价或局限 \\
    \midrule
    流水线与模板 & 需要可控、可逐步检查的生成 & 内容确定、文档规划、微规划、表层实现由规则或模板完成 & 规则按领域人工编写，组合随字段增多而膨胀，表达单一 \\
    端到端编码器—解码器 & 去掉人工规则，从数据学出选择与措辞 & 记录嵌入、注意力与门控复制，抽取式评价衡量忠实性 & 无依据内容多，内容选择不稳定，所有决定隐含在逐词生成中 \\
    内容选择与规划 & 选择与次序缺少全局约束 & 内容选择门、指针网络规划、计划编码后再生成 & 计划标签靠弱对齐；多一次规划解码；漏选无法补回 \\
    线性化与预训练模型 & 平行数据少，表达能力受限 & 结构展开为带分隔符的序列，微调T5或BART，或提示大语言模型 & 对顺序敏感，长表需截断，内容选择重新变为隐式 \\
    \bottomrule
  \end{tabularx}
  \caption{数据到文本的方法谱系。自上而下，前一行的局限大致是后一行的动机；后一行并不在所有维度上都优于前一行，例如模板的事实覆盖最容易保证。}
  \label{tab:nlp-generation-d2t-methods}
\end{table*}

\section{开放写作}
\label{sec:nlp-generation-open-writing}

开放写作包括故事、续写和带风格要求的自由文本。它允许虚构，内容不再由输入记录逐条决定，挑战随之改变：长文本要前后连贯、情节要有组织；采样需要在多样性与第~\ref{sec:nlp-inference-degeneration}小节的重复退化之间取得平衡；常常还要控制情感、话题或正式程度等属性；长篇中的人物设定和已写事实必须保持一致。

方法谱系按这些挑战排列。分层故事生成先产生计划或提示，再据此写正文，处理长程连贯。属性控制有三条路线：训练时把属性作为条件，解码时用属性模型的梯度修改隐状态，解码时用属性分类器对候选重加权；第一条要训练整个生成器，后两条冻结生成器、只外接小模型，第三条还省去了解码时的反向传播。数千词的长篇超出单次生成的连贯范围，需要递归规划与修订。大语言模型直接续写再对候选重排，是所有路线的对照基线。节末表~\ref{tab:nlp-generation-writing-methods}汇总比较。

本节固定一个写作任务。题目为“失而复得的借书卡”，人物只有林岚，故事发生在同一天的图书馆，分三段，每段一句，文体平实；允许虚构找卡经过，但不得新增人物替林岚找到卡，也不得把图书馆业务规则写成真实规定。为小说构造情节不属于事实错误，前提是任务明确允许虚构；评价分别检查情节多样性、阅读偏好与约束满足，任何一个维度都不能代表全部创作质量。

\subsection{分层故事生成}
\label{sec:nlp-generation-hierarchical-story}

单层语言模型从标题直接写长故事，常在几句之后偏离主题：每一步只依据前文局部地决定下一个词，开头的线索没有在结尾回收，远处的句子之间缺少共同的约束。\term{分层生成}（hierarchical generation）先产生一个高层计划，再以计划为条件写正文，使相距很远的句子共享同一条件。两种代表路线的区别在于计划从哪里来：Plan-and-Write从故事中自动抽取关键词作为计划监督，Fan等的分层生成以人写的写作提示作为上层。

\subsubsection{Plan-and-Write情节规划}
\label{sec:nlp-generation-plan-write}

Plan-and-Write先用有序的关键词表示故事线，再依据故事线生成文本\scite{nlp-generation-yao2019}。它在五句短故事语料上训练，用RAKE关键词抽取算法从每句抽出一个关键词，形成与句子一一对应的故事线，作为规划监督。故事线只是关键词序列，不是完备的事件逻辑，也不是人物状态记录。

原方法比较两种规划方式。静态规划由题目$u$一次生成完整故事线$z_{1:k}$，再生成故事$y$，两个模型分别最大化$p(z\mid u)$与$p(y\mid u,z)$。动态规划交替进行：以题目、此前句子和上一关键词预测下一关键词，再据此写下一句，规划条件是$(u,y_{<j},z_{j-1})$，写作条件还包含新关键词$z_j$。两个模型可以分别训练、运行时串接。静态方式在写正文之前就决定了全部情节，便于检查首尾呼应；动态方式能响应已写内容，但计划和正文交替推进，前面的线索可能在后面被遗忘。

静态教学计划为“遗失／回找／夹页”，据此形成三段：
\begin{quote}
林岚在图书馆借书时发现借书卡不见了，记起自己刚翻过一本厚书。

她返回阅览桌，把那本厚书从原处取下来，一页页查看。

借书卡正夹在书页之间；她收好卡，回到借书处。
\end{quote}
另设一份教学检查记录：第一段林岚在借书处，卡的位置未知；第二段她到阅览桌，取书后开始找；第三段卡从书页转到她手中，她返回借书处。开头的厚书线索在结尾回收，人物移动和事件先后可以逐项检查。这份记录是本书的核验手段，不是原模型维护的状态。

动态方式写完第一段后，也可能选择“衣袋”作为第二个关键词，形成“她先检查外套衣袋，找到夹在票据中的借书卡”，再以“收好”结尾。这时第一段的“厚书”线索失去作用；若第三段仍生成“从书中取出刚找到的卡”，同一张卡就出现在两个位置。静态规划也会失败：计划写“衣袋”而正文仍沿用“夹页”时，只有核对最终文本才能发现计划与实现脱节。

\subsubsection{提示到故事的分层生成}
\label{sec:nlp-generation-prompt-story}

Fan等把计划换成一句人写的写作提示，建立了提示—故事配对的WritingPrompts数据集，数据来自网络写作社区，故事平均长达数百词\scite{nlp-inference-fan2018story}。生成分两层：先用卷积语言模型生成一句提示，再以提示为条件生成故事。提示比关键词序列信息更丰富，也更接近人类写作时的起点。

故事模型是卷积编码器—解码器。卷积只能看到有限窗口，故事又长达数百词，他们在解码器中加入门控多尺度自注意力，让每个位置能读取全部已写内容：查询、键和值不用线性投影，而由带门控线性单元的小网络给出；各头在不同时间尺度上读取历史，第一个头看全部位置，第二个头每隔一步取一个位置，第三个头每隔两步取一个，依此类推，使不同的头关注不同范围的信息。更关键的问题是，在提示—故事数据上直接训练的编码器—解码器常常忽视提示，只学会写通顺的故事：词与词之间的局部依赖比故事与提示之间的依赖容易学得多。他们的对策是\term{模型融合}（model fusion）：先训练一个编码器—解码器并固定参数，再训练第二个编码器—解码器，第二个模型每步读取第一个模型的隐状态。设两者在第$t$步的隐状态为$\bm h^{\mathrm{new}}_t$与$\bm h^{\mathrm{pre}}_t$，融合层先拼接再逐维门控，
\[
  \bm g_t=\sigma\bigl(\bm W[\bm h^{\mathrm{new}}_t;\bm h^{\mathrm{pre}}_t]+\bm b\bigr),\qquad
  \bm h_t=\bm g_t\odot[\bm h^{\mathrm{new}}_t;\bm h^{\mathrm{pre}}_t],
\]
再经若干全连接层进入输出层。Fan等把这看作一种残差学习：第一个模型已经承担了大部分语言建模，第二个模型可以集中学习第一个模型没有学到的部分，例如故事对提示的依赖。

生成时，Fan等使用top-k采样（$k=10$），每步只在概率最高的十个词元中重新归一化抽样；束搜索在开放写作中产生重复、平淡的文本，原因及截断采样的一般形式见第~\ref{sec:nlp-generation-sampling}小节。与Plan-and-Write相比，提示—故事路线的上层更自然，但提示本身不规定情节结构，长故事中段仍可能离题；两者都只做一层规划，第~\ref{sec:nlp-generation-long-form}小节把规划扩展到多层。

\subsection{属性控制生成}
\label{sec:nlp-generation-attribute-control}

许多写作任务还要求文本具有指定属性$a$：情感为正面、话题为科技、措辞正式。目标是从条件分布$p(x\mid a)$生成。由贝叶斯公式，$p(x\mid a)\propto p(x)\,p(a\mid x)$；逐词元写开，第$t$步应满足
\begin{equation}
  p(x_t\mid x_{<t},a)\propto p(x_t\mid x_{<t})\,p(a\mid x_{\le t}),
  \label{eq:nlp-generation-attribute-bayes}
\end{equation}
其中$p(a\mid x_{\le t})$是“当前前缀写完后满足属性”的概率。三条路线以不同方式实现这个关系。CTRL直接训练条件模型$p(x\mid a)$，控制强但新增属性需重训；PPLM冻结语言模型，用属性模型的梯度修改隐状态，近似地把分布推向高$p(a\mid x)$的区域，不需重训，但每步要多次反向传播；FUDGE训练一个前缀分类器估计$p(a\mid x_{\le t})$，按式~\eqref{eq:nlp-generation-attribute-bayes}对候选重加权，只需前向计算。表~\ref{tab:nlp-generation-attribute-control}在本小节末比较三者。

\subsubsection{控制码条件生成：CTRL}
\label{sec:nlp-generation-ctrl}

CTRL是一个约16亿参数的Transformer语言模型\scite{nlp-inference-keskar2019ctrl}。训练时在每条序列开头放一个\term{控制码}（control code），即标明这条文本来源或类型的特殊前缀，学习条件语言模型
\[
  p(x\mid a)=\prod_{t}p(x_t\mid x_{<t},a).
\]
控制码来自语料的来源信息：维基百科、图书、评论等领域，评论的评分，网页的网址，以及问答、翻译等任务标记。生成时写下“Reviews Rating: 5.0”再接开头，模型就按五星评论的风格续写。

控制码把属性控制变成普通的条件生成，推断不增加任何开销，控制效果也最直接。它的局限来自训练：控制码集合在训练时就固定了，只能控制语料中天然带有标签的属性；想要“正式而简短”这类新属性或任意组合，只能重新训练或微调。CTRL还提出了第~\ref{chap:nlp-model-inference}章介绍的重复惩罚。

\subsubsection{梯度引导控制：PPLM}
\label{sec:nlp-generation-pplm}

\term{即插即用语言模型}（plug and play language model, PPLM）的出发点是不改动已经训练好的语言模型，只在解码时外接一个小的属性模型\scite{nlp-generation-dathathri2020pplm}。属性模型给出$p(a\mid x)$，有两种：词袋模型取一组属性词，以下一步分布落在这组词上的总概率作为$p(a\mid x)$；判别器模型在语言模型隐状态的平均上接一个线性分类器。属性模型的参数很少，可以按需更换。

控制通过修改语言模型的历史状态实现。设$\bm H_t$是第$t$步之前各层的键值缓存，PPLM给它加一个扰动$\Delta\bm H_t$，沿属性对数概率的梯度方向更新若干次：
\begin{equation}
  \Delta\bm H_t\leftarrow\Delta\bm H_t+\eta\,
  \frac{\nabla_{\Delta\bm H_t}\log p(a\mid\bm H_t+\Delta\bm H_t)}
  {\bigl\lVert\nabla_{\Delta\bm H_t}\log p(a\mid\bm H_t+\Delta\bm H_t)\bigr\rVert^{\gamma}},
  \label{eq:nlp-generation-pplm}
\end{equation}
$\eta$是步长，分母按梯度范数的$\gamma$次方归一化，控制每次修改的幅度。只追求属性会损害流畅性，因此更新目标中还加了一项KL惩罚，使修改后的下一词元分布$\widetilde p_{t+1}$不远离原分布$p_{t+1}$；采样时再取二者的几何加权平均$\widetilde p_{t+1}^{\,\gamma_{\mathrm{gm}}}p_{t+1}^{\,1-\gamma_{\mathrm{gm}}}$并重新归一化，$\gamma_{\mathrm{gm}}$接近1时更偏向控制后的分布。

\Needspace{12\baselineskip}
\begin{algorithm}[PPLM的单步生成]
  \label{alg:nlp-generation-pplm}
  \begin{algorithmic}[1]
    \Require 冻结的语言模型；历史键值$\bm H_t$；属性模型$p(a\mid\cdot)$；步数$m$，步长$\eta$，系数$\gamma,\lambda_{\mathrm{KL}},\gamma_{\mathrm{gm}}$
    \Ensure 下一词元$x_{t+1}$与更新后的历史
    \State 用$\bm H_t$前向一次，得到原分布$p_{t+1}$；令$\Delta\bm H_t\gets\bm 0$
    \For{$i=1,\ldots,m$}
      \State 用$\bm H_t+\Delta\bm H_t$前向，得到$\widetilde p_{t+1}$与$\log p(a\mid\bm H_t+\Delta\bm H_t)$
      \State 目标取$\log p(a\mid\cdot)-\lambda_{\mathrm{KL}}\operatorname{KL}(\widetilde p_{t+1}\Vert p_{t+1})$，反向传播到$\Delta\bm H_t$
      \State 按式~\eqref{eq:nlp-generation-pplm}更新$\Delta\bm H_t$
    \EndFor
    \State 以$\widetilde p_{t+1}^{\,\gamma_{\mathrm{gm}}}p_{t+1}^{\,1-\gamma_{\mathrm{gm}}}$归一化后采样$x_{t+1}$
    \State 用修改后的历史前向$x_{t+1}$，追加其键值，进入下一步
  \end{algorithmic}
\end{algorithm}

算法~\ref{alg:nlp-generation-pplm}的每一步都要做$m$次前向与反向传播，生成速度比普通采样慢数倍；步长过大时文本会变得不通顺甚至重复属性词，需要调节$\eta,\lambda_{\mathrm{KL}},\gamma_{\mathrm{gm}}$。原文每步更新$m=3$到$10$次，$\lambda_{\mathrm{KL}}=0.01$，$\gamma_{\mathrm{gm}}$取$0.8$到$0.95$，梯度范数的归一化对Transformer每层分别计算。与CTRL相比，PPLM不需重训且属性模型可任意更换；与下文的FUDGE相比，它修改的是隐状态而不是输出概率，代价是每步的反向传播。

\subsubsection{FUDGE软属性控制}
\label{sec:nlp-generation-fudge}

正式、礼貌等属性难以写成一个词袋。FUDGE直接实现式~\eqref{eq:nlp-generation-attribute-bayes}：训练一个前缀分类器$q_{\bm\phi}(a\mid x_{\le t})$，预测“当前前缀继续写完后，完整文本会不会满足属性”，再用它对下一词元重加权\scite{yang-klein-2021-fudge}。训练样本是带完整文本属性标签的序列，把同一个完整文本的标签附给它的每个前缀，用二元交叉熵训练；分类器判断的是前缀的未来走向，不是前缀本身是否已经正式。生成器保持不变。

记生成器的前缀为$h=x_{<t}$，候选下一词元为$v$。重加权后的分布为
\begin{equation}
  \widetilde p(v\mid h,a)=
  \frac{p_{\bm{\theta}}(v\mid h)\,q_{\bm{\phi}}(a\mid hv)}
  {\sum_w p_{\bm{\theta}}(w\mid h)\,q_{\bm{\phi}}(a\mid hw)}.
  \label{eq:nlp-generation-fudge}
\end{equation}
输入材料也属于条件，这里为简洁略去。分母要对每个候选调用一次分类器，实际实现只对基础分布概率最高的若干词元计算，这也改变了可搜索的集合。分类器只是近似，不一定与生成器的真实未来分布校准一致，因此公式实现的是控制的近似，不是属性满足的保证。

给定开头“校外来宾”，三个教学候选片段为“须提交”“记得交”“无需提交”，基础概率为$(0.40,0.50,0.10)$，预测最终正式的概率为$(0.90,0.20,0.95)$。乘积为$(0.36,0.10,0.095)$，总和$0.555$，归一化为$(0.649,0.180,0.171)$。束宽为一时，选择从口语的“记得交”变成“须提交”。“无需提交”同样很正式，控制后它的相对概率反而从$0.10$升到$0.171$，却与通知要求矛盾：软属性控制只管措辞，内容是否正确要单独核验。

FUDGE只需前向计算，分类器可以比生成器小得多，多个属性的分类器还可以相乘组合；后者只有在相应条件独立假设下才有概率解释，“极短”与“完整保留长条件”这类属性还可能直接冲突。与PPLM相比，FUDGE不修改隐状态，控制强度受限于基础分布前若干候选中有没有满足属性的词元。

\begin{table*}[htbp]
  \centering
  \small
  \begin{tabularx}{\linewidth}{@{}>{\raggedright\arraybackslash}p{18mm}
    >{\raggedright\arraybackslash}X>{\raggedright\arraybackslash}X
    >{\raggedright\arraybackslash}X>{\raggedright\arraybackslash}X@{}}
    \toprule
    方法 & 控制的实现 & 训练代价 & 推断代价 & 控制强度与流畅性 \\
    \midrule
    CTRL & 训练条件模型$p(x\mid a)$，控制码放在序列开头 & 整个模型在带来源标签的语料上训练；新属性需重训 & 与普通生成相同 & 控制直接；属性集合固定，难以组合 \\
    PPLM & 沿$\nabla\log p(a\mid x)$修改键值历史，KL项与几何平均保持流畅 & 只训练小属性模型 & 每步$m$次前向与反向 & 控制较强；步长过大损害流畅性 \\
    FUDGE & 前缀分类器估计$p(a\mid x_{\le t})$，对候选按式~\eqref{eq:nlp-generation-fudge}重加权 & 训练前缀分类器 & 每步对前若干候选各做一次分类器前向 & 不改动生成器，流畅性好；强度受候选集合限制 \\
    \bottomrule
  \end{tabularx}
  \caption{三种属性控制方法的比较。三者都实现式~\eqref{eq:nlp-generation-attribute-bayes}的某种近似：CTRL在训练时学习条件分布，PPLM和FUDGE在解码时引入外部属性模型。属性控制只改变措辞或话题倾向，不检查内容是否与条件一致。}
  \label{tab:nlp-generation-attribute-control}
\end{table*}

\subsection{长篇规划与递归修订}
\label{sec:nlp-generation-long-form}

前两小节的方法处理的是几百词的故事。数千词的长篇超出了单次生成能保持连贯的范围：上下文窗口装不下全部已写内容，即使装得下，模型也会逐渐偏离最初的设定，前文定下的人物属性在后文被改写。分层故事生成只有一层计划，不足以约束这么长的文本。

Re3（recursive reprompting and revision）用大语言模型的提示把长篇写作拆成“计划—草拟—重排—编辑”的循环\scite{nlp-generation-yang2022re3}。计划阶段由前提生成背景、人物及其描述和分条大纲。草拟阶段每次写一段，把与当前大纲条目相关的计划内容、前文摘要和最近写出的几段重新组织进提示，生成多个候选续写；这种每段重建提示的做法称为\term{递归重提示}（recursive reprompting），使远处的设定以摘要形式始终留在上下文中。重排阶段用两个训练好的打分模型选出与大纲最相关、与前文最连贯的候选。编辑阶段从已写内容中抽取人物属性，检测新段落与之矛盾的地方，再调用编辑模型修改。

\Needspace{12\baselineskip}
\begin{algorithm}[递归重提示与修订的长篇生成]
  \label{alg:nlp-generation-re3}
  \begin{algorithmic}[1]
    \Require 前提$u$；大语言模型；相关性与连贯性打分模型；每段候选数$N$；段数上限$M$
    \Ensure 长篇文本$y$
    \State 由$u$生成背景、人物描述与大纲条目$o_1,\ldots,o_K$；令$y\gets\varnothing$
    \For{$j=1,\ldots,M$}
      \State 选出当前大纲条目$o_k$，组织提示：相关计划、前文摘要、最近若干段
      \State 生成$N$条候选段落
      \State 按与$o_k$的相关性和与$y$的连贯性重排，取最高者$y_j$
      \State 抽取人物属性，检测$y_j$与已写内容的矛盾，必要时编辑$y_j$
      \State 把$y_j$接到$y$后，更新前文摘要；当前条目写完则转到$o_{k+1}$
    \EndFor
    \State \Return $y$
  \end{algorithmic}
\end{algorithm}

DOC在Re3的基础上加深计划\scite{nlp-generation-yang2023doc}。它把大纲逐层细化成多级条目，每个条目附有场景和出场人物，细化时对候选子条目过滤和重排；生成每一段时，用一个按FUDGE方式实现的控制器（基于OPT-350m）逐词元调整候选的概率，使正文贴合当前的细化条目。控制器用段落与其摘要的配对做对比训练，负例包括流畅但与摘要无关的续写，使长段落不只是开头切题。这里第~\ref{sec:nlp-generation-fudge}小节的属性从“正式”换成了“与某个大纲条目相关”。

长篇规划路线的代价有三项。其一，调用次数随段数和候选数成倍增加，一篇长篇需要数百次模型调用和多个辅助打分模型。其二，错误沿层级传播：大纲层的矛盾会被忠实地写进每一段，编辑阶段只能修补局部，无法改正结构问题。其三，评价长篇的连贯和有趣主要依赖人工判断，自动指标难以反映，第~\ref{sec:nlp-generation-evaluation}节将讨论这一问题。

\subsection{语言模型直接续写与候选重排}
\label{sec:nlp-generation-continuation}

作为所有路线的对照基线，可以把题目、人物要求和已写段落作为前缀，让大语言模型直接续写，生成方式与采样规则沿用第~\ref{sec:nlp-generation-sampling}小节。运行时要写清模型能否看到此前所有段落、用何种段落标记，以及第三段结束后的EOS；没有EOS而撞上长度上限的文本仍是截断候选。多轮共同修改题目或人物设定的交互过程见第~\ref{chap:nlp-question-answering-dialogue}章。

直接续写的结果需要筛选和重排，方法见第~\ref{sec:nlp-inference-reranking}小节。固定上一小节故事的第一段，教学候选甲续写“返回阅览桌找书，卡夹在书中”，候选乙续写“管理员把卡送来，林岚感谢他”。乙新增人物代为找卡，违反已定条件，即使语言偏好评分更高也先被淘汰。在合法候选之间，再按线索回收、重复和约定风格重排；重排器可以是人工量表、规则或有标注的偏好模型，来源不同需要分别验证。

比较规划路线与直接续写时，要把规划消耗的词元与调用计入预算。例如设两条完整候选、每条最多90个新词元：规划路线每条分配12个规划词元和78个正文词元，直接路线每条90个正文词元，两者再接受相同次数的核验，记录实际词元、前向调用和延迟（教学协议）。规划在写正文之前决定部分情节，直接续写在正文中逐步决定；完整上下文可用时二者都能读取已写内容，发生截断时都可能丢失人物条件。偏好最高的一条仍须检查位置、时间线、题目要求和允许虚构的边界。

表~\ref{tab:nlp-generation-writing-methods}汇总本节方法。规划、控制和修订各自回应开放写作的一项挑战，可以组合使用；DOC就是分层规划与FUDGE式控制的组合。

\begin{table*}[htbp]
  \centering
  \small
  \begin{tabularx}{\linewidth}{@{}>{\raggedright\arraybackslash}p{24mm}
    >{\raggedright\arraybackslash}X>{\raggedright\arraybackslash}X
    >{\raggedright\arraybackslash}X@{}}
    \toprule
    方法 & 解决的问题 & 核心机制 & 新代价或局限 \\
    \midrule
    Plan-and-Write & 单层生成离题，首尾不呼应 & 自动抽取的关键词故事线，静态或动态规划后写作 & 关键词不表达事件逻辑；计划与正文可能脱节 \\
    提示到故事 & 计划信息太少；模型忽视条件 & 先生成提示，卷积编码器—解码器加多尺度自注意力，模型融合 & 只有一层规划，长故事中段仍会离题 \\
    CTRL & 需要指定风格或领域 & 训练时以控制码为条件 & 属性集合固定，新属性需重训 \\
    PPLM & 不重训而外接属性 & 属性梯度修改键值历史 & 每步多次反向传播，流畅性受步长影响 \\
    FUDGE & 降低外接控制的代价 & 前缀分类器对候选重加权 & 近似未校准；强度受候选集合限制 \\
    Re3与DOC & 数千词长篇的连贯与人物一致 & 多层大纲，递归重提示，候选重排与一致性编辑 & 调用次数多；大纲错误沿层级传播；评价依赖人工 \\
    直接续写与重排 & 作为对照基线，实现最简单 & 前缀续写，截断采样，规则或偏好模型重排 & 无显式计划；长文本中设定易被遗忘 \\
    \bottomrule
  \end{tabularx}
  \caption{开放写作的方法谱系。前两行处理长程连贯，中间三行处理属性控制，Re3与DOC把规划扩展到长篇；直接续写是各路线的共同基线。比较时须计入规划、控制与重排消耗的调用和词元。}
  \label{tab:nlp-generation-writing-methods}
\end{table*}
